\section*{Sources and draft status}
The SHA-256 word functions, constants, initial state, block processing and padding follow
\href{https://doi.org/10.6028/NIST.FIPS.180-4}{NIST FIPS 180-4, Secure Hash Standard},
sections 4 through 6. The official record's August 2015 version is used; the technical
SHA-256 definition is unchanged from its preceding edition.
\href{https://www.rfc-editor.org/rfc/rfc6234}{RFC 6234} provides a second primary
description, test material and an implementation reference. The discussion of keyed
hashing distinguishes the primitive from the construction specified in
\href{https://www.rfc-editor.org/rfc/rfc2104}{RFC 2104}. The functions, context layout,
workspace ownership, limits and status values in this chapter are local teaching contracts.

The private prototypes compare complete digests with JDK SHA-256. A separate arithmetic
model checks schedules and compression from arbitrary state words; a Python calculation
checks the worked values against complete-message digests from hashlib. A register-specific
first-round check applies only to the worked illustration. Equivalent solutions are
assessed by observable results and ownership at the public function interfaces.

This is a directly edited native chapter foundation. Its cumulative student edits,
step-specific quizzes and hints, observation recipes and distribution archives are still
being developed. The measured work table reports verified function-boundary instruction
counts. It is not yet included as the twenty-fifth package
in the main manuscript. Real annotated IDE captures, remote lesson delivery, enforced
progression and independent accessibility validation remain pending. The author-requested
substantial rewrite for length, continuity and explanatory quality remains deferred until
all thirty-five chapter drafts exist.
